\documentclass[10pt,a4paper]{amsart}
\usepackage{amsmath,amssymb,amsthm,graphicx}
% Retained minimal preamble for the Stacks Project chapter coherent.tex (Stacks chapter, pre-delegation retained fragment for stacks-coherent).
% Verbatim excerpt of the chapter preamble (sources/stacks-spaces-descent/preamble.tex): the reference/slogan/history
% environment block (native lines 7-12), Xy-pic setup (native lines 14-19) and the native theorem declarations with
% the mathematical macros (native lines 154-217). Copied mechanically; nothing added, removed or redefined. The
% wrapper class and the amsmath/amssymb/amsthm/graphicx base packages are supplied by the assembler, so the
% class/xr-hyper/hyperref/externaldocument part of the native preamble is omitted.
% For dealing with references we use the comment environment
\usepackage{verbatim}
\newenvironment{reference}{\comment}{\endcomment}
%\newenvironment{reference}{}{}
\newenvironment{slogan}{\comment}{\endcomment}
\newenvironment{history}{\comment}{\endcomment}
% For commutative diagrams we use Xy-pic
\usepackage[all]{xy}

% We use 2cell for 2-commutative diagrams.
\xyoption{2cell}
\UseAllTwocells
% Theorem environments.
%
\theoremstyle{plain}
\newtheorem{theorem}[subsection]{Theorem}
\newtheorem{proposition}[subsection]{Proposition}
\newtheorem{lemma}[subsection]{Lemma}

\theoremstyle{definition}
\newtheorem{definition}[subsection]{Definition}
\newtheorem{example}[subsection]{Example}
\newtheorem{exercise}[subsection]{Exercise}
\newtheorem{situation}[subsection]{Situation}

\theoremstyle{remark}
\newtheorem{remark}[subsection]{Remark}
\newtheorem{remarks}[subsection]{Remarks}

\numberwithin{equation}{subsection}

% Macros
%
\def\lim{\mathop{\mathrm{lim}}\nolimits}
\def\colim{\mathop{\mathrm{colim}}\nolimits}
\def\Spec{\mathop{\mathrm{Spec}}}
\def\Hom{\mathop{\mathrm{Hom}}\nolimits}
\def\Ext{\mathop{\mathrm{Ext}}\nolimits}
\def\SheafHom{\mathop{\mathcal{H}\!\mathit{om}}\nolimits}
\def\SheafExt{\mathop{\mathcal{E}\!\mathit{xt}}\nolimits}
\def\Sch{\mathit{Sch}}
\def\Mor{\mathop{\mathrm{Mor}}\nolimits}
\def\Ob{\mathop{\mathrm{Ob}}\nolimits}
\def\Sh{\mathop{\mathit{Sh}}\nolimits}
\def\NL{\mathop{N\!L}\nolimits}
\def\CH{\mathop{\mathrm{CH}}\nolimits}
\def\proetale{{pro\text{-}\acute{e}tale}}
\def\etale{{\acute{e}tale}}
\def\QCoh{\mathit{QCoh}}
\def\Ker{\mathop{\mathrm{Ker}}}
\def\Im{\mathop{\mathrm{Im}}}
\def\Coker{\mathop{\mathrm{Coker}}}
\def\Coim{\mathop{\mathrm{Coim}}}

% Boxtimes
%
\DeclareMathSymbol{\boxtimes}{\mathbin}{AMSa}{"02}

%
% Macros for moduli stacks/spaces
%
\def\QCohstack{\mathcal{QC}\!\mathit{oh}}
\def\Cohstack{\mathcal{C}\!\mathit{oh}}
\def\Spacesstack{\mathcal{S}\!\mathit{paces}}
\def\Quotfunctor{\mathrm{Quot}}
\def\Hilbfunctor{\mathrm{Hilb}}
\def\Curvesstack{\mathcal{C}\!\mathit{urves}}
\def\Polarizedstack{\mathcal{P}\!\mathit{olarized}}
\def\Complexesstack{\mathcal{C}\!\mathit{omplexes}}
% \Pic is the operator that assigns to X its picard group, usage \Pic(X)
% \Picardstack_{X/B} denotes the Picard stack of X over B
% \Picardfunctor_{X/B} denotes the Picard functor of X over B
\def\Pic{\mathop{\mathrm{Pic}}\nolimits}
\def\Picardstack{\mathcal{P}\!\mathit{ic}}
\def\Picardfunctor{\mathrm{Pic}}
\def\Deformationcategory{\mathcal{D}\!\mathit{ef}}
\usepackage{cleveref}
\crefname{originallemma}{original lemma}{original lemmas}
\Crefname{originallemma}{Original lemma}{Original lemmas}
\makeatletter
\newlabel{algebra-lemma-Artin-Rees}{{00IN}{source}{}{sourcepointer-1}{}}
\newlabel{algebra-lemma-Artin-Rees@cref}{{[originallemma][0][]00IN}{[1][0][]source}{}{}{}}
\makeatother
\crefname{originallemma}{original lemma}{original lemmas}
\Crefname{originallemma}{Original lemma}{Original lemmas}
\makeatletter
\newlabel{algebra-lemma-completion-Noetherian-Noetherian}{{0316}{source}{}{sourcepointer-2}{}}
\newlabel{algebra-lemma-completion-Noetherian-Noetherian@cref}{{[originallemma][0][]0316}{[1][0][]source}{}{}{}}
\makeatother
\crefname{originallemma}{original lemma}{original lemmas}
\Crefname{originallemma}{Original lemma}{Original lemmas}
\makeatletter
\newlabel{algebra-lemma-map-AR}{{00IO}{source}{}{sourcepointer-3}{}}
\newlabel{algebra-lemma-map-AR@cref}{{[originallemma][0][]00IO}{[1][0][]source}{}{}{}}
\makeatother
\crefname{originallemma}{original lemma}{original lemmas}
\Crefname{originallemma}{Original lemma}{Original lemmas}
\makeatletter
\newlabel{lemma-inverse-systems-affine}{{087W}{source}{}{sourcepointer-4}{}}
\newlabel{lemma-inverse-systems-affine@cref}{{[originallemma][0][]087W}{[1][0][]source}{}{}{}}
\makeatother
\title{The category of coherent formal modules is abelian: for X Noetherian and I a quasi-coherent ideal sheaf, Coh(X, I) is abelian, restriction to open subsets is exact, and exactness can be checked on the members of an open covering.}
\author{The Stacks Project Authors}
\date{stacks-project:coherent.tex:local-snapshot-20261012; GNU FDL 1.2 or later}
\begin{document}
\maketitle
\noindent\textit{Editorial note.} Statement span is lines 6529-6541 of sources/stacks-coherent/pinned-original-source.tex; proof span is lines 6543-6610, the complete author proof (66 body lines).\par\medskip
\noindent\textit{Editorial note.} Navigation stubs: 4 cited-but-not-included labels are supplied as non-load-bearing pointers with their pinned tags where the snapshot tag index has them.\par\medskip
\noindent\textit{Editorial note.} extra\_packages is empty on purpose: the retained preamble already loads xy with options and a second declaration clashes in the assembler wrapper.\par\medskip
\noindent\textit{Editorial note.} Text is retained verbatim, including the author's own typos and placeholders.\par\medskip
\noindent\textit{Editorial note.} Disclosure per the tightened ruling: the termwise cokernel claim carries '(details omitted)' and parts (2)-(3) are called immediate from the construction; the entire kernel construction and its Artin-Rees analysis - the load-bearing content - is written in the span.\par\medskip
\noindent\textit{Editorial note.} Span computed from the \textbackslash{}label marker through the enclosing \textbackslash{}end\{lemma\}; precard\_gate RESULT=OK (probe L9999).\par\medskip
\noindent\textit{Editorial note.} Cross-chapter same-label warning (NOT an occupancy conflict, Chief ruling 2026-10-13 evening): label 'lemma-inverse-systems-abelian' also exists in stacks-spaces-more-morphisms — a different proposition in a different chapter file.\par\medskip
\noindent\textit{Allowed standard prerequisites (editorial).}\begin{itemize}
\item Explicit prerequisite boundary. See why-field; every pointer named here occurs verbatim in the retained statement/proof text.
\end{itemize}\medskip
\ifdefined\hypertarget\hypertarget{sourcepointer-1}{}\fi
\noindent\textit{Source reference.} Original lemma 00IN: sources/stacks-coherent/pinned-original-source.tex, source lines 6610--6610. Cited by the retained statement or proof but not part of the collected spans; supplied only as a navigation pointer. The argument does not depend on its content.\par\medskip
\ifdefined\hypertarget\hypertarget{sourcepointer-2}{}\fi
\noindent\textit{Source reference.} Original lemma 0316: sources/stacks-coherent/pinned-original-source.tex, source lines 6610--6610. Cited by the retained statement or proof but not part of the collected spans; supplied only as a navigation pointer. The argument does not depend on its content.\par\medskip
\ifdefined\hypertarget\hypertarget{sourcepointer-3}{}\fi
\noindent\textit{Source reference.} Original lemma 00IO: sources/stacks-coherent/pinned-original-source.tex, source lines 6610--6610. Cited by the retained statement or proof but not part of the collected spans; supplied only as a navigation pointer. The argument does not depend on its content.\par\medskip
\ifdefined\hypertarget\hypertarget{sourcepointer-4}{}\fi
\noindent\textit{Source reference.} Original lemma 087W: sources/stacks-coherent/pinned-original-source.tex, source lines 6489--6495. Cited by the retained statement or proof but not part of the collected spans; supplied only as a navigation pointer. The argument does not depend on its content.\par\medskip
\begin{lemma}
\label{lemma-inverse-systems-abelian}
Let $X$ be a Noetherian scheme and let $\mathcal{I} \subset \mathcal{O}_X$
be a quasi-coherent sheaf of ideals.
\begin{enumerate}
\item The category $\textit{Coh}(X, \mathcal{I})$ is abelian.
\item For $U \subset X$ open the restriction functor
$\textit{Coh}(X, \mathcal{I}) \to \textit{Coh}(U, \mathcal{I}|_U)$
is exact.
\item Exactness in $\textit{Coh}(X, \mathcal{I})$ may be checked by
restricting to the members of an open covering of $X$.
\end{enumerate}
\end{lemma}

\begin{proof}
Let $\alpha =(\alpha_n) : (\mathcal{F}_n) \to (\mathcal{G}_n)$ be a morphism of
$\textit{Coh}(X, \mathcal{I})$. The cokernel of $\alpha$ is the inverse system
$(\Coker(\alpha_n))$ (details omitted). To describe the kernel let
$$
\mathcal{K}'_{l, m} = \Im(\Ker(\alpha_l) \to \mathcal{F}_m)
$$
for $l \geq m$.
We claim:
\begin{enumerate}
\item[(a)] the inverse system $(\mathcal{K}'_{l, m})_{l \geq m}$ is
eventually constant, say with value $\mathcal{K}'_m$,
\item[(b)] the system $(\mathcal{K}'_m/\mathcal{I}^n\mathcal{K}'_m)_{m \geq n}$
is eventually constant, say with value $\mathcal{K}_n$,
\item[(c)] the system $(\mathcal{K}_n)$ forms an object of
$\textit{Coh}(X, \mathcal{I})$, and
\item[(d)] this object is the kernel of $\alpha$.
\end{enumerate}
To see (a), (b), and (c) we may work affine locally, say $X = \Spec(A)$
and $\mathcal{I}$ corresponds to the ideal $I \subset A$. By
Lemma \ref{lemma-inverse-systems-affine}
$\alpha$ corresponds to a map $f : M \to N$ of finite $A^\wedge$-modules.
Denote $K = \Ker(f)$. Note that $A^\wedge$ is a Noetherian
ring (Algebra, Lemma \ref{algebra-lemma-completion-Noetherian-Noetherian}).
Choose an integer $c \geq 0$ such that
$K \cap I^n M \subset I^{n - c}K$ for $n \geq c$
(Algebra, Lemma \ref{algebra-lemma-Artin-Rees})
and which satisfies Algebra, Lemma \ref{algebra-lemma-map-AR}
for the map $f$ and the ideal $I^\wedge = IA^\wedge$. Then
$\mathcal{K}'_{l, m}$ corresponds to the $A$-module
$$
K'_{l, m} = \frac{a^{-1}(I^lN) + I^mM}{I^mM} =
\frac{K + I^{l - c}f^{-1}(I^cN) + I^mM}{I^mM} =
\frac{K + I^mM}{I^mM}
$$
where the last equality holds if $l \geq m + c$. So $\mathcal{K}'_m$
corresponds to the $A$-module $K/K \cap I^mM$ and
$\mathcal{K}'_m/\mathcal{I}^n\mathcal{K}'_m$ corresponds to
$$
\frac{K}{K \cap I^mM + I^nK} = \frac{K}{I^nK}
$$
for $m \geq n + c$ by our choice of $c$ above. Hence $\mathcal{K}_n$
corresponds to $K/I^nK$.

\medskip\noindent
We prove (d). It is clear from the description on affines above that
the composition $(\mathcal{K}_n) \to (\mathcal{F}_n) \to (\mathcal{G}_n)$
is zero. Let $\beta : (\mathcal{H}_n) \to (\mathcal{F}_n)$
be a morphism such that $\alpha \circ \beta = 0$. Then
$\mathcal{H}_l \to \mathcal{F}_l$ maps into $\Ker(\alpha_l)$.
Since $\mathcal{H}_m = \mathcal{H}_l/\mathcal{I}^m\mathcal{H}_l$
for $l \geq m$ we obtain a system of maps
$\mathcal{H}_m \to \mathcal{K}'_{l, m}$. Thus a map
$\mathcal{H}_m \to \mathcal{K}_m'$. Since
$\mathcal{H}_n = \mathcal{H}_m/\mathcal{I}^n\mathcal{H}_m$ we obtain
a system of maps $\mathcal{H}_n \to \mathcal{K}'_m/\mathcal{I}^n\mathcal{K}'_m$
and hence a map $\mathcal{H}_n \to \mathcal{K}_n$ as desired.

\medskip\noindent
To finish the proof of (1) we still have to show that $\Coim = \Im$
in $\textit{Coh}(X, \mathcal{I})$. We have seen above that taking
kernels and cokernels commutes, over affines, with the description
of $\textit{Coh}(X, \mathcal{I})$ as a category of modules. Since
$\Im = \Coim$ holds in the category of modules
this gives $\Coim = \Im$ in $\textit{Coh}(X, \mathcal{I})$.
Parts (2) and (3) of the lemma are immediate from our construction
of kernels and cokernels.
\end{proof}

\end{document}
